\documentclass[12pt]{article}
\usepackage[utf8]{inputenc}
\usepackage[T1]{fontenc}
\usepackage{hyperref}
\title{Yan-nhangu Sign Language Language Resources for K-12 Curriculum Development}
\author{Compiled by Claude Sonnet 5.5 (claude-sonnet-5-5) for Andr\'as Kornai}
\date{}
\begin{document}
\maketitle

\section{Introduction}
Yan-nhangu Sign Language (ISO 639-3 yhs) is an endangered shared sign language of the Crocodile Islands off Arnhem Land, Australia. It is used as a first language by the elderly only. The search results mainly describe the wider Yolngu Sign Language (a separate code, ygs), so what is specific to Yan-nhangu is uncertain. No Yan-nhangu-specific materials were found.

\section{Primary Educational Resources}
No primers, textbooks or school programmes were found.

\section{Linguistic and Reference Materials}
\begin{itemize}
\item Nhangu language (Wikipedia), on the spoken language associated with the sign language. \url{https://en.wikipedia.org/wiki/Nhangu_language}
\item Yolngu sign language (PH Museum photographic project). It concerns the wider Yolngu sign language. \url{https://phmuseum.com/projects/yolngu-sign-language?f=f}
\item Academic record on Yolngu sign language (University of Central Lancashire repository). Its exact content was not confirmed from the snippet. \url{https://clok.uclan.ac.uk/14387}
\end{itemize}
The search summary mentioned an illustrated handbook of over 500 signs documented by linguist Bentley James. No URL for it was seen, so it is not listed as a resource.

\section{Religious and Community Resources}
None found.

\section{Digital and Audio Resources}
None specific to Yan-nhangu were found.

\section{Curriculum Development Context}
Signing is used in ceremony, hunting and mourning, and children traditionally learn it by everyday observation. Any curriculum would be a revitalization effort led by Yolngu signers, with community permission, building on the broader Yolngu Sign Language documentation. Further searching is needed.

\end{document}